\section{Uniqueness of the continuous prefix}\label{sec:collisions}

We now pass from representations to distinct totient values.  We retain the
tuple sets $\mathcal T(t)$, $\mathcal G(t)$ and $\mathcal D(t)$ of
Section~\ref{sec:extraction}: the first consists of the distinct tuples
$(p_0,\ldots,p_R,d)$ admitting a basic witness, the second consists of those
for which every basic witness satisfies \eqref{eq:good-conditions}, and
$\mathcal D(t)=\mathcal T(t)\setminus\mathcal G(t)$.  Their value map is
\[
 F(p_0,\ldots,p_R,d)=d\prod_{j=0}^R(p_j-1).
\]
The tail preimages of $d$ are not part of a tuple.  This distinction is
essential: we shall prove that almost every value has one such tuple,
while allowing several tail preimages.

\subsection{Counting pairs with the same value}
The alignment of largest factors and cancellation follow the argument
around (5.34)--(5.37) in Ford's Section~5 \cite{FordTotients2013}.
Here we count all ordered pairs of distinct tuples, including the data
inside both residual factors.
We first cancel the common prefix and align the large factors in the
remaining shifts. After applying the comparison estimate, we recover
the tuple coordinates inside the residual factors and then restore
the common prefix.

\begin{proposition}[Prefix collisions]\label{prop:collisions}
Fix $c>1$, and form all parameters from $x$, with $H$ fixed before $x$
tends to infinity.  For $t=x$ or $t=x/c$, let
\[
 K(t)=\#\{(\tau,\sigma)\in\mathcal G(t)^2:
             \tau\ne\sigma,\ F(\tau)=F(\sigma)\}.
\]
There is a quantity $\epsilon_H\to0$, independent of the phase, such that
\[
 K(t)\le\bigl(\epsilon_H+o_{x;H}(1)\bigr)
                   \frac{xG_m}{\log x}.
\]
The pairs counted here are ordered pairs of distinct tuples.
\end{proposition}

\begin{proof}
Choose one basic witness for each tuple under consideration.  All choices
satisfy \eqref{eq:good-conditions}, by the definition of $\mathcal G(t)$.
If two distinct tuples have the same value, some displayed prime differs:
equality of all the primes $p_0,\ldots,p_R$ would also force equality of
$d$.  Let $i$ be the first unequal prime index, and write $q_j$ for the
primes in the other witness.  Put
\[
 h=m-i,\qquad J=\lceil30\log h\rceil,\qquad
 z=z_i=\exp\exp(0.7b_{i+J}),\qquad S=S_i=\exp\exp(b_i^{1/3}).
\]
Thus $h\ge H$.  The residual integers in the good conditions are
\[
 D_i=\varphi(p_{i+J+1}\cdots p_La),\qquad
 D_i'=\varphi(q_{i+J+1}\cdots q_La').
\]
Canceling the common prefix leaves the identity
\begin{equation}\label{eq:suffix-identity}
 v^{(i)}=D_i\prod_{j=0}^{J}(p_{i+j}-1)
        =D_i'\prod_{j=0}^{J}(q_{i+j}-1).
\end{equation}
We will apply Proposition~\ref{prop:comparison} to this identity.  First
we verify its size and interval hypotheses; then we count all ways in
which its solutions can give tuple pairs.

\paragraph{Sizes of the suffix and its layers.}
For sufficiently large $H$, uniformly for $h\ge H$,
\begin{equation}\label{eq:collision-scales}
 b_i\asymp h\rho^{-h},\qquad
 \log b_i=O(h),\qquad
 c_1h^{-20}\le\frac{b_{i+J}}{b_i}
       =(1-J/h)\rho^J\le C_1h^{-18}.
\end{equation}
Indeed $30\lambda$ lies strictly between $18$ and $20$; the ceiling in
$J$ changes only constant factors.  In particular $J<h/2$ and
$i+J<L$.  The constants in this proof are independent of $H,x,i$ and the
phase, once $H$ and then $x$ are sufficiently large.

For $i=0$ set $y=x$.  For $i\ge1$ place $v^{(i)}$ in a dyadic interval
$(y/2,y]$.  The coarse bands in \eqref{eq:basic} imply that the logarithm
of the preimage suffix after $p_i$ is at most
\[
 h\exp(0.62b_i)+\exp(2\widetilde b_L).
\]
To see the first term, use
$1.1\widetilde b_{i+1}<0.62b_i$ for large $x$, and bound the number of
displayed factors by $h$.  The bound for $a$ gives the second term.
These quantities are negligible compared with
$\log p_i\ge\exp(0.88b_i)$.  Consequently, writing $B_y=\log_2y$,
\begin{equation}\label{eq:suffix-sizes}
 p_i,q_i\ge y^{0.9},\qquad
 0.88b_i\le B_y\le1.12b_i,\qquad u_i\le B_y+O(1).
\end{equation}
For $i=0$ the first inequality is part of the basic witness conditions,
and $B_y=b_0=u_0=B$ by convention.

Our target is a bound of the form
\[
 \frac{y}{\log y}\exp(-\kappa B_y/h^4)
\]
for the suffix pairs, with an absolute $\kappa>0$.
The strict simplex slack will supply the exponential saving.
Since $B_y\asymp b_i$, this saving can absorb the
$\exp(O(b_i\log h/h^8))$ cost of recovering tuple coordinates
inside $D_i$ and $D_i'$.

Let
\[
 \delta=\sqrt{\log_2S/B_y}\asymp b_i^{-1/3}.
\]
It follows from \eqref{eq:collision-scales} that $S<z$ and
$J\delta=o(h^{-20})$.  Normality gives
$\Omega(p_{i+j}-1),\Omega(q_{i+j}-1)\le4B_y$ for $0\le j\le J$.
For example, split the prime factors at $S$, apply the two normality
conditions, and use $p_{i+j}-1\le v^{(i)}\le y$.  If
\[
 A_j=P^+(p_{i+j}-1),\qquad A_j'=P^+(q_{i+j}-1),
\]
then
\begin{equation}\label{eq:largest-factor-band}
 \log_2A_j=u_{i+j}+O(\log b_i),\qquad
 \log_2A_j'=\log_2q_{i+j}+O(\log b_i).
\end{equation}
This follows from
$(p_{i+j}-1)\le A_j^{\Omega(p_{i+j}-1)}$ and the reverse bound
$A_j\le p_{i+j}-1$.  When $i=j=0$, the convention $u_0=B$
causes only an $O(1)$ difference from $\log_2p_0$, since
$x^{0.9}\le p_0\le x+1$.

The bands in \eqref{eq:largest-factor-band} are disjoint and decrease
with $j$.  Their smallest separation is a constant multiple of
$b_i h^{-20}$, which dominates the error $O(\log b_i)$.  All the $A_j,A_j'$
exceed $z$, whereas
\[
 P^+(D_i),P^+(D_i')<z.
\]
For the latter assertion, every prime in the residual preimage is at
most $p_{i+J+1}$ or $q_{i+J+1}$, and
$1.1\widetilde b_{i+J+1}<0.7b_{i+J}$; taking a totient cannot increase
the largest prime factor.  Repetitions of the last prime in $a$ or $a'$
do not affect this conclusion.

\paragraph{Aligning the largest factors on the two sides.}
We claim that
\begin{equation}\label{eq:factor-alignment}
 |\log_2 A_j-\log_2 A_j'|\le(2j+1)\delta B_y
                 \qquad(0\le j\le J).
\end{equation}
Suppose, for example, that $A_j>A_j'$, and put
$e=(\log_2 A_j-\log_2 A_j')/B_y$.  On the first side of
\eqref{eq:suffix-identity}, each of its first $j+1$ shifts contributes
at least $(e-\delta)B_y$ prime factors, with multiplicity, to
$(A_j',A_j]$.  On the other side only its first $j$ shifts can
contribute, each at most $(e+\delta)B_y$.  The later shifts have smaller
largest factors, and both residual integers are $z$-smooth.  The lower
endpoint exceeds $S$, so normality applies.  Equality of the two sides
therefore gives
\[
 (j+1)(e-\delta)\le j(e+\delta),
\]
which proves \eqref{eq:factor-alignment}.  Interchanging the sides
handles the other order.

Cancel any equal primes among the two lists in
\eqref{eq:suffix-identity}, and let $r$ be the product of their shifts.
The disjoint prime bands show that equality is possible only at the
same original index.  The index $i$ survives, since $p_i\ne q_i$.
Let $b\le J+1\le B_y$ be the length of the remaining lists.  Their corresponding
primes are unequal, as required by Proposition~\ref{prop:comparison}.

Set $Y_b=z$, $Y_0=y$ and $\log_2U_0=0.8B_y$.  At each remaining
positive index, choose an interval on a $\delta$-grid in normalized
double logarithms which contains the corresponding pair from
\eqref{eq:factor-alignment}.  Its width is $O(J\delta)$.  The intervals
remain strictly separated, because $J\delta=o(h^{-20})$.  Each lies
below height $0.8$, by \eqref{eq:suffix-sizes} and the upper bound
$u_{i+1}\le0.62b_i$.  At index zero,
\eqref{eq:suffix-sizes} and \eqref{eq:largest-factor-band} give largest
factors above height $0.9$.  Thus all the ordered-cutoff hypotheses of
Proposition~\ref{prop:comparison} hold.  If $b=1$, there are no positive
indices to grid and $Y_1=z$.

The part of $p_i-1$ supported on primes at most $Y_1$ has logarithm
at most $4B_y\log Y_1\le4B_y\exp(0.8B_y)$.  Since $p_i\ge y^{0.9}$,
the complementary part exceeds $\sqrt y$ for large $H$.  The whole
common product is squarefree above $Y_b$, by the good conditions;
canceling the equal shifts preserves this property.  This verifies
the last two hypotheses of Proposition~\ref{prop:comparison}.

Apply that proposition to the remaining lists with $D=D_i$, $D'=D_i'$
and the factor $r$ just canceled; their common product is at most $y/r$.
Write $\nu_k=\log_2Y_k/B_y$ and $\mu_k=\log_2U_k/B_y$ for the
normalized interval endpoints.

Write $j_1<\cdots<j_{b-1}$ for the surviving positive original
indices.  Since $j_k\ge k$ and $a_k\le a_{j_k}$, the grid intervals
and the strict slack condition in \eqref{eq:good-conditions} give
\begin{align}
 \sum_{k=1}^{b-1}a_k\nu_k
 &\le \frac{1}{B_y}\sum_{j=1}^{J}a_j u_{i+j}
                          +O(J^2\log(2J)\delta)\notag\\
 &\le1-h^{-4}+O(B_y^{-1}+J^2\log(2J)\delta).
 \label{eq:collision-slack}
\end{align}
Here $a_j$ is increasing and $\sum_{j\le J}a_j=O(J\log(2J))$.
The grid widths and the explicit $\delta$ term in
\eqref{eq:comparison} contribute another
$O(J^2\log(2J)\delta)=o(h^{-4})$ to its exponent of $\log y$.
Thus, for some absolute $\kappa>0$, its logarithmic factor is at most
\[
 (\log y)^{-1-\kappa h^{-4}}
 =\frac1{\log y}\exp(-\kappa B_y/h^4).
\]
It remains to show that the other factors and the recovery of tuple
coordinates use only a smaller part of this saving.

\paragraph{Recovering tuples and summing the remaining choices.}
The comparison estimate counts its displayed primes and its residual
integer $D_i'$.  We must also account for the tuple data forgotten
inside $D_i,D_i'$.  If $R\le i+J$, all displayed tuple primes are
already determined.  Their tail totient is
\[
 d=D_i\prod_{j=R+1}^{i+J}(p_j-1),
\]
and similarly on the other side, with an empty product when $R=i+J$.
If $R>i+J$, the missing tuple data give an ordered factorization
\begin{equation}\label{eq:residual-factorization}
 D_i=d\prod_{j=i+J+1}^{R}(p_j-1).
\end{equation}
There are at most $(h+1)^{\Omega(D_i)}$ such factorizations.  Indeed,
for any integer $n$, assigning its $\Omega(n)$ prime occurrences to
$k$ labeled slots bounds the number of ordered multiplicative
factorizations into $k$ factors by $k^{\Omega(n)}$.  Once a factor
$p_j-1$ is fixed, $p_j$ is fixed as well.  This argument also covers
repeated prime factors.  Consequently the two missing tails cost at most
\begin{equation}\label{eq:tuple-recovery-cost}
 (h+1)^{\Omega(D_i)+\Omega(D_i')}
       \le\exp\{2b_i\log(h+1)/h^8\}.
\end{equation}
The factorizations in \eqref{eq:residual-factorization} are exact:
the primes up to $R$ are distinct and the remaining preimage tail is
smaller than $p_R$, by \eqref{eq:cofactor-size}.  All repetitions inside
that tail are kept in its totient $d$.

We now sum the comparison bound over the canceled primes, the grids,
and $D_i$.  The canceled index sets contribute at most $2^{J+1}$.
Each canceled prime has reciprocal shift sum $O(b_i)$ by Mertens'
theorem.  The number of grid choices is at most
$(C/\delta)^{2(J+1)}$.  Thus their total cost, including the factors
$(CB_y^6)^b$ in \eqref{eq:comparison}, is $\exp(O(h^2))$.
For the residual integer, the finite Euler product gives
\[
 \sum_{P^+(D)\le z}\frac1D\ll\log z.
\]
Before enlarging this sum we use the good-condition bounds
$\Omega(D_i),\Omega(D_i')\le b_i/h^8$ in both
\eqref{eq:tuple-recovery-cost} and the factor $(b+1)^{\Omega(D_i)}$
of the comparison estimate.

The logarithm of all the remaining costs, apart from the factor
$y/(\log y)$ and the saving from \eqref{eq:collision-slack}, is at most
\begin{equation}\label{eq:collision-costs}
 O(h^2)+O\left(\frac{b_i\log h}{h^8}\right)
       +O\left(\frac{b_i\log h\log\log(3h)}{h^{18}}\right)
       =o\left(\frac{b_i}{h^4}\right).
\end{equation}
The last term uses
$\log_2z\ll b_i h^{-18}$ to bound the factor
$(\log Y_b)^{1+3b\log(b(b+1))}$ in \eqref{eq:comparison}, as well
as the sum over $D_i$.  All little-oh estimates here are uniform for
$h\ge H$ as $H\to\infty$, because $b_i\asymp h\rho^{-h}$.
We have therefore obtained, for some absolute $c_2>0$,
\begin{equation}\label{eq:suffix-pair-count}
 \#\{\text{ordered suffix tuple pairs with }y/2<v^{(i)}\le y\}
 \ll\frac{y}{\log y}\exp\{-c_2B_y/h^4\}.
\end{equation}
For $i=0$ the same expression bounds all the pairs with first unequal
index zero, using $y=x$.  Choosing one witness for each tuple has not
introduced multiplicity into the desired count: comparison data and
the factorizations just counted determine each ordered pair of tuples.

\paragraph{Summing suffixes and restoring the common prefix.}
For $i\ge1$ the required sum over suffix pairs is weighted by
$1/v^{(i)}$.  A dyadic interval in \eqref{eq:suffix-pair-count}
contributes at most
\[
 \frac{C}{\log y}\exp\{-c_2\log_2y/h^4\}.
\]
Writing the dyadic endpoints as $y=2^n$, and comparing the sum with
an integral, gives
\begin{align}
 \sum_{\text{suffix pairs}}\frac1{v^{(i)}}
 &\ll\sum_{\log(n\log2)\ge0.88b_i-O(1)}
       \frac{\exp\{-c_2\log(n\log2)/h^4\}}{n\log2}\notag\\
 &\ll\int_{0.88b_i-O(1)}^\infty\exp(-c_2u/h^4)\,\dd u
  \ll h^4\exp(-c_3b_i/h^4).
 \label{eq:dyadic-suffix-mass}
\end{align}
In this change of variable, $u=\log(n\log2)$, the factor
$1/\log y$ exactly compensates for the density of dyadic endpoints.
In particular, no factor of order $\exp(b_i)$ is lost.

For fixed suffix data and common primes $p_1,\ldots,p_{i-1}$, the
number of possible $p_0$ is at most
\[
 \frac{Cx}{\log x}\,
 \frac1{v^{(i)}\prod_{j=1}^{i-1}(p_j-1)}.
\]
Indeed the upper endpoint for $p_0$ is
$1+x/(v^{(i)}\prod_{j=1}^{i-1}(p_j-1))$; if the range is nonempty,
its restriction $p_0\ge x^{0.9}$ makes its logarithm comparable to
$\log x$, so the prime number theorem gives this upper bound.
The projected-prefix estimate \eqref{eq:projected-mass} and
\eqref{eq:dyadic-suffix-mass} now show that the contribution for this
$i$ is at most
\[
 \frac{CxG_m}{\log x}\,h^4\exp(-c_3b_i/h^4).
\]
This also applies to the empty projected prefix when $i=1$.
Summing over $h\ge H$ yields the convergent tail
\[
 C\sum_{h\ge H}h^4\exp\{-c_4\rho^{-h}/h^3\}\longrightarrow0.
\]
For $i=0$, \eqref{eq:suffix-pair-count} gives
$O((x/\log x)\exp\{-c_2B/m^4\})$, which is
$o(xG_m/\log x)$.  This proves the proposition, uniformly for both
endpoints $t$.
\end{proof}

\subsection{From tuples to values and least preimages}

The following elementary lemma keeps track of exceptional tuples as
well as exceptional values.  This is needed because a small set of
values could otherwise carry many representations.

\begin{lemma}[Finite-map counting]\label{lem:finite-map}
Let $T$ and $W$ be finite sets, let $F:T\to W$, and write
$T=G\sqcup D$.  Suppose that every element of $W\setminus E$ is
in $F(T)$, where $E\subset W$.  For $w\in W$ put
$r_w=|G\cap F^{-1}(w)|$, and let
\[
 K=\sum_{w\in W}r_w(r_w-1).
\]
Then
\[
 -|E|\le |T|-|W|\le |D|+K/2.
\]
Moreover, with
$E^*=E\cup F(D)\cup\{w:r_w\ge2\}$, one has
\[
 |E^*|\le|E|+|D|+K/2,\qquad
 |F^{-1}(E^*)|\le|E|+2|D|+K.
\]
The restriction of $F$ to $T\setminus F^{-1}(E^*)$ is a bijection
onto $W\setminus E^*$.
\end{lemma}

\begin{proof}
Coverage gives the lower bound.  For every positive integer $r$,
$r-1\le r(r-1)/2$, and hence
\[
 |G|-|F(G)|\le K/2.
\]
Adding the at most $|D|$ remaining elements proves the upper bound.
There are at most $K/2$ fibers with $r_w\ge2$, which proves the
bound for $|E^*|$.  Also
\[
 \sum_{w\in E^*}r_w\le|E^*|+K/2,
\]
by $r\le1+r(r-1)/2$ for integers $r\ge0$.  Adding the elements of
$D$ proves the second bound.  Outside $E^*$ coverage, absence of
discarded elements, and $r_w\le1$ give exactly one preimage under $F$.
\end{proof}

\begin{proposition}[A common prefix for all preimages]\label{prop:unique-prefix}
For $t=x$ or $t=x/c$, with the same parameters formed from $x$,
\begin{equation}\label{eq:tuple-cardinality}
 \bigl|\,|\mathcal T(t)|-V(t)\,\bigr|
       \le\bigl(\epsilon_H+o_{x;H}(1)\bigr)\frac{xG_m}{\log x},
 \qquad\epsilon_H\longrightarrow0.
\end{equation}
There is a set $E^*(t)\subset\mathcal V\cap[1,t]$ for which both
$|E^*(t)|$ and the number of tuples in $\mathcal T(t)$ with values in
$E^*(t)$ are at most
$(\epsilon_H+o_{x;H}(1))xG_m/\log x$, after enlarging $\epsilon_H$.
For every $v\in(\mathcal V\cap[1,t])\setminus E^*(t)$ there is a unique tuple
$(p_0,\ldots,p_R,d)\in\mathcal T(t)$ of value $v$, and every preimage
of $v$ has this prefix and a remaining factor of totient $d$.
In particular,
\begin{equation}\label{eq:least-factorization}
 \ell(v)=p_0\cdots p_R\ell(d).
\end{equation}
\end{proposition}

\begin{proof}
Apply Lemma~\ref{lem:finite-map} with $T=\mathcal T(t)$,
$G=\mathcal G(t)$, $D=\mathcal D(t)$ and
$W=\mathcal V\cap[1,t]$.  Let $E$ be the exceptional values in
Proposition~\ref{prop:basic}, restricted to $[1,t]$.  Every preimage
of each $v\notin E$ is basic, so it induces an element of $T$.
The size of $E$ is negligible on the scale $xG_m/\log x$, by
\eqref{eq:ford-scale}; Proposition~\ref{prop:discard} gives the same
bound for $|D|$, and Proposition~\ref{prop:collisions} gives it for
$K$.  The counting conclusions therefore follow from the lemma.

Fix $v\notin E^*(t)$.  Every preimage is basic and so induces a tuple
in $T$.  It cannot induce a tuple in $D$, since $v\notin F(D)$.
The remaining tuple is unique.  Thus, writing $Q=p_0\cdots p_R$,
every preimage of $v$ has the form $Qz$, where $\varphi(z)=d$.
It follows that every such preimage is at least $Q\ell(d)$.

For the reverse inequality choose one basic witness with tail $w$.
The deterministic size bound \eqref{eq:cofactor-size} gives
$\ell(d)\le w<p_R$.  The prefix primes are distinct, so
$\gcd(Q,\ell(d))=1$ and
\[
 \varphi(Q\ell(d))=\varphi(Q)\varphi(\ell(d))
                 =d\prod_{j=0}^R(p_j-1)=v.
\]
Hence $\ell(v)\le Q\ell(d)$, proving \eqref{eq:least-factorization}.
\end{proof}

The universal quantifier on preimages in Proposition~\ref{prop:basic}
is used only after the counting argument, to obtain the common-prefix
statement and \eqref{eq:least-factorization}.  A single basic preimage
would give the upper bound for $\ell(v)$ but would not give the reverse
inequality.
